Un long solénoïde de rayon $r_s = \qty{\v{rs}}{cm}$ compte \num{\v{n}} spires par mètre. Le courant qui le parcourt augmente à un rythme constant de \qty{\v{dIdt}}{A/s}. Une spire circulaire de rayon $r_\ell = \qty{\v{rl}}{cm}$ et de résistance \qty{\v{R}}{\ohm} entoure le solénoïde ; elle est coaxiale avec lui et placée loin de ses extrémités. ($\mu_0 = 4\pi \times \qty{e-7}{T.m/A}$)

\begin{enumerate}[label=\alph*)]
    \item Déterminez le taux de variation du champ magnétique à l'intérieur du solénoïde.
    \item Déterminez la force électromotrice induite dans la spire.
    \item Déterminez le courant induit dans la spire.
    \item Le courant induit circule-t-il dans le même sens que le courant du solénoïde ou en sens inverse ? Justifiez.
\end{enumerate}
